\documentclass[../main.tex]{subfiles}

\begin{document}
\section{\textbf{Classical Statistical Mechanics}}
\makebox[\textwidth][s]{Author: Xianchao \hfill Date: \today}

In this section, we will derive the probability distribution of different ensemble. Here ``ensemble'' means the system under different specific boundary conditions. Typically, there are  microcanonical ensemble, canonical ensemble, grand canonical ensemble, and $NPT$ ensemble. Microcanonical ensemble is the system that has the number of particles $N$, the total volume $V$, and the total energy $E$, fixed. Canonical ensemble is the system that has $N$, $V$ fixed, while allowing exchange energy with surroundings. Grand canonical ensemble allows system to exchange not only energy, but the particles. Only requires $V$ fixed. For $NPT$ ensemble, it's obviously that, it fixes the $N$, temperature $T$ and pressure $P$. In general, the probability distribution for all kinds of ensemble, can be derived by maximizing the Shannon entropy, with constraints applied.

\subsection{Microcanonical Ensemble}

Let's look into the easy case, the microcanonical ensemble. In microcanonical ensemble, The total number of particles $N$, total volume $V$, and total energy $E$, are fixed. Assume allowed specific configuration $\sigma_i$ will appear at probability $p_i$, then the only constraints we can apply on the system is the normalization condition, which is:

\begin{equation}
    \sum_{i} p_i = 1
\end{equation}

The Lagrange multiplies can be written as:

\begin{equation}
    L = -\sum_{i} p_i \ln p_i - \alpha(\sum_{i} p_i - 1)
\end{equation}

Taking derivative gives the equation:

\begin{equation}
    \pdv{L}{p_i} = -1 - \ln p_i - \alpha = 0
\end{equation}

We can immediately obtain $p_i = \textrm{const}$. Assuming there's $W$ allowed configurations for given $N$, $V$, $E$, then the probability of any one specific configuration is:

\begin{equation}
    p_i = \frac{1}{W}
\end{equation}

And the entropy is:

\begin{equation}
    S = -k_B \sum_{i} p_i \ln p_i = k_B \ln W
\end{equation}

\subsection{Canonical Ensemble}

For the canonical ensemble, which fixes the particle number $N$, the volume $V$. The total energy for a specific configuration $\sigma_i$ is no longer fixed, but we know the average energy $U$:

\begin{equation}
    U = \sum_{i} E_i p_i = \langle E \rangle
\end{equation}

Also, there is the probability normalization condition, then the Lagrange multiplies is:

\begin{equation}
    L = -\sum_{i} p_i \ln p_i - \alpha(\sum_{i} p_i - 1) - \beta(\sum_{i} E_i p_i - U)
\end{equation}

Taking the derivative gives the equation:

\begin{equation}
    \pdv{L}{p_i} = -1 - \ln p_i - \alpha - \beta E_i = 0
\end{equation}

Which gives:

\begin{equation}
    p_i \propto \textrm{e}^{-\beta E_i}
\end{equation}

We can set the normalization factor as $\frac{1}{Z}$, where $Z$ is the partition function. Now we can get that:

\begin{equation}
    p_i = \frac{\textrm{e}^{-\beta E_i}}{Z}
\end{equation}

With:

\begin{equation}
    Z = \sum_{i} \textrm{e}^{-\beta E_i}
\end{equation}

Be careful, now the parameter $\beta$ is just a Lagrange multiplier. And we haven't proven that it is the signal that commonly used in physics, $\frac{1}{k_B T}$. To show this, we have to build the connection with thermodynamics. Insert the expression of probability into the information entropy, we can get:

\begin{align}
    \frac{S}{k_B} &= -\sum_{i} p_i (-\beta E_i - \ln Z)\nonumber\\
    &= \beta \sum_{i} p_i E_i + \ln Z\nonumber\\
    &= \beta U + \ln Z
\end{align}

So that we have:

\begin{equation}
    \frac{1}{k_B}\dd{S} = \beta \dd{U} + U\dd{\beta} + \pdv{\ln Z}{\beta}\dd{\beta}
\end{equation}

Also:

\begin{align}
    \pdv{\ln Z}{\beta} = \frac{1}{Z}(-\sum_{i} \textrm{e}^{-\beta E_i} E_i) = -U
\end{align}

Combine this two equations, we have:

\begin{equation}
    \dd{S} = k_B\beta \dd{U}
\end{equation}

Which gives $\pdv{S}{U}|_{N, V} = k_B\beta$. For the canonical ensemble, the equilibrium state satisfies the deferential relation:

\begin{equation}
    \dd{U} = T\dd{S} - p\dd{V}
\end{equation}

Considering the volume is fixed, so $\dd{V} = 0$, which gives:

\begin{equation}
    \dd{U} = T\dd{S}
\end{equation}

It naturally gives that $\pdv{S}{U}|_{N, V} = \frac{1}{T}$. So we have:

\begin{equation}
    \beta = \frac{1}{k_B T}
\end{equation}

\subsection{Grand Canonical Ensemble}

For the grand canonical ensemble, both energy and particle number are allowed to change. Only the volume $V$ is fixed. In this case, we have constraints:

\begin{align}
    \bar{N} &= \sum_{i} p_i N_i =\langle N \rangle \\
    U &= \sum_{i} E_i p_i = \langle E \rangle
\end{align}

In this case, we can write the Lagrange multiplier as:

\begin{equation}
    L = -\sum_{i} p_i \ln p_i - \alpha (\sum_{i} p_i - 1) - \beta (\sum_{i} E_i p_i - U) - \gamma (\sum_{i} p_i N_i - \bar{N})
\end{equation}

Taking derivative and we have:

\begin{equation}
    \pdv{L}{p_i} = -1 - \ln p_i - \alpha - \beta E_i - \gamma N_i = 0
\end{equation}

In this case we have:

\begin{equation}
    p_i \propto \textrm{e}^{-\beta E_i - \gamma N_i}
\end{equation}

And we can use the same approach to normalize the probability distribution by partition function:

\begin{align}
    p_i &= \frac{1}{Z}\textrm{e}^{-\beta E_i - \gamma N_i}\\
    \Xi &= \sum_{i} \textrm{e}^{-\beta E_i - \gamma N_i}
\end{align}

Now we need to solve out the Lagrange multiplier by matching the expression to the thermodynamics. According to the basic thermodynamics relation, we have:

\begin{equation}
    \dd{U} = T\dd{S} - p \dd{V} + \mu\dd{N}
\end{equation}

The volume is fixed in this ensemble, which can give:

\begin{align}
    \dd{S} &= \frac{p}{T}\dd{V} - \frac{\mu}{T}\dd{N} + \frac{1}{T}\dd{U}\nonumber\\
    &= - \frac{\mu}{T}\dd{N} + \frac{1}{T}\dd{U}
\end{align}

The expression for Shannon entropy is:

\begin{align}
    S &= -k_B \sum_{i} p_i \ln p_i\nonumber\\
    &= -k_B \sum_{i} p_i (-\beta E_i - \gamma N_i - \ln \Xi)\nonumber\\
    &= k_B \beta U + k_B\gamma \bar{N} + k_B \ln \Xi
\end{align}

Which gives:

\begin{equation}
    \dd{S} = k_B \beta \dd{U} + k_B U \dd{\beta} + k_B\bar{N}\dd{\gamma} + k_B\gamma\dd{N} + k_B \dd{\ln \Xi}
\end{equation}

Also, we have:

\begin{align}
    \dd{\ln \Xi} &= \frac{1}{\Xi}(\sum_{i} -E_i\textrm{e}^{-\beta E_i - \gamma N_i}\dd{\beta} - N_i\textrm{e}^{-\beta E_i - \gamma N_i}\dd{\gamma})\nonumber\\
    &= -U\dd{\beta} - \bar{N}\dd{\gamma}
\end{align}

By which we have:

\begin{equation}
    \dd{S} = k_B \beta \dd{U} + k_B\gamma \dd{N}
\end{equation}

By matching this expression with that from thermodynamics, we have:

\begin{align}
    \beta = \frac{1}{k_B T}\\
    \gamma = - \frac{\mu}{k_B T}
\end{align}

Till now, we successively obtain the expression for the grand canonical ensemble:

\begin{equation}
    p_i = \frac{1}{\Xi} \textrm{e}^{-\frac{E_i - \mu N_i}{k_B T}}
\end{equation}

\subsection{$NPT$ Ensemble}

The $NPT$ ensemble works in the same pattern. The number of particles, pressure, temperature are fixed, while the energy $E$ and volume $V$ is allowed to change. In this case we have:

\begin{align}
    \sum_{i} p_i V_i &= \bar{V}\\
    \sum_{i} p_i E_i &= U
\end{align}

Using the same approach, we can get:

\begin{equation}
    p_i \propto \textrm{e}^{-\beta E_i - \gamma V_i}
\end{equation}

The partition function is:

\begin{equation}
    Z = \sum_{i} \textrm{e}^{-\beta E_i - \gamma V_i}
\end{equation}

And the thermodynamics can give out the relation:

\begin{equation}
    \dd{S} = \frac{1}{T}\dd{U} + \frac{P}{T}\dd{V}
\end{equation}

And we can get the expression of entropy as:

\begin{equation}
    \frac{S}{k_B} = \beta U + \gamma \bar{V} + \ln Z
\end{equation}

Take derivative on $\ln Z$:

\begin{align}
    \dd{\ln Z} &= (\sum_{i}-E_i p_i \dd{\beta} - \sum_{i}V_i p_i \dd{\gamma})\nonumber\\
    &= -U\dd{\beta} - \bar{V}\dd{\gamma}
\end{align}

By matching the coefficients, we have:

\begin{align}
    k_B \beta = \frac{1}{T}\\
    k_B \gamma = \frac{P}{T}
\end{align}

Thus, we can get:

\begin{equation}
    p_i = \frac{\textrm{e}^{-\frac{E_i + PV_i}{k_B T}}}{Z}
\end{equation}



\end{document}
